Groups scan in parallel and merge computed-primary survivors \cite{johnson_gpu}.

\subsection{Candidate partition and deterministic ordering}

For a fixed query, concatenate the selected inverted lists in probe order
to obtain $n_c$ candidate occurrences. Assume that each visited database
row occurs once and that the index stays fixed during search. For
$\gscan\in\{1,2,4,8\}$, block $(q,j)$ owns the interval
\begin{equation}
 U_j=\left[\left\lfloor\frac{j n_c}{\gscan}\right\rfloor,
              \left\lfloor\frac{(j+1)n_c}{\gscan}\right\rfloor\right),
 \qquad j=0,\ldots,\gscan-1.
 \label{eq:5-partition}
\end{equation}
Here an interval denotes the corresponding candidates, rather than their
database IDs. The device computes its endpoints from the routed candidate
count. Groups differ in size by at most one candidate; a long list can
cross a group boundary. This balances candidate counts, whereas assigning
equal numbers of lists need not do so. Each thread finds its first list
through the probe-prefix offsets and advances its list cursor while
scanning its interval. Residual codes are not read at this stage.

The selection key is
\begin{equation}
 K_q(x)=\bigl(\primarydist(q,x),p(x)\bigr),
 \qquad
 (a,p)\prec(b,r)\iff a<b\ \text{or}\ (a=b\ \text{and}\ p<r),
 \label{eq:5-primary-key}
\end{equation}
where $p(x)$ is the physical row position in the fixed inverted-list
storage. It is unique among the query's candidates and independent of
the order of atomic appends. It is not the database ID used to break
final-result ties. With non-NaN scores, this defines a total order on
distinct candidates. Local selection and merging use this order. Partition independence also requires identical computed scores across paths.

\subsection{Exact local retention and global coverage}

$\operatorname{Top}_c(V)$ contains the first $\min(c,|V|)$ keys, where $c=\max\{k,\lceil\alpha k\rceil\}$. Each group keeps up to $c$, since all global winners may lie there; a $c/\gscan$ quota is unsafe.

\begin{resultblock}
\begin{lemma}[Local-buffer retention]
\label{lem:5-buffer}
Let $V_t$ be the candidates scored by a block through chunk $t$, and
$B_t$ its valid buffer after any compaction or filtering at that point.
Assume fixed complete keys, sufficient append capacity and exact
top-$c$ compaction. Let $\beta_t$ be the minimum of the complete shared
bounds observed so far, including the initial top sentinel. Then
\begin{equation}
 \operatorname{Top}_c(V_t)\cap\{x:K_q(x)\preceq\beta_t\}
 \subseteq B_t.
 \label{eq:5-buffer-invariant}
\end{equation}
With sharing disabled, the bound is the sentinel and the buffer contains
the whole prefix top-$c$.
\end{lemma}
\begin{proof}
The empty prefix is immediate. Put
\[
 E_t=\operatorname{Top}_c(V_t)\cap\{x:K_q(x)\preceq\beta_t\}.
\]
On extending $V_t$ to $V_{t+1}$, $\beta_{t+1}\preceq\beta_t$ and
\[
 \operatorname{Top}_c(V_{t+1})\cap V_t\subseteq\operatorname{Top}_c(V_t).
\]
Hence every old member of $E_{t+1}$ was retained. A locally rejected new key has $c$ earlier predecessors, so cannot enter the new prefix top-$c$. All other eligible arrivals fit by the capacity assumption. For any eligible winner $x$,
\[
 \#\{y\in B:K_q(y)\prec K_q(x)\}
 \leq\#\{y\in V_{t+1}:K_q(y)\prec K_q(x)\}<c,
\]
so exact compaction keeps $x$. Filtering removes only keys above an observed bound; their running minimum covers stale reads. These steps preserve the invariant inductively.
\end{proof}
\end{resultblock}

\begin{resultblock}
\begin{proposition}[Coverage and truncated merging]
\label{prop:5-local-coverage}
For fixed keys under a common total order, let $U$ be the visited
universe and $V$ any retained subset. If
$\operatorname{Top}_c(U)\subseteq V\subseteq U$, then
$\operatorname{Top}_c(V)=\operatorname{Top}_c(U)$.
Pairwise sorted merging with truncation to $c$ preserves this result.
In particular, for disjoint groups,
\begin{equation}
 \operatorname{Top}_c\!\left(\bigcup_j U_j\right)
 =\operatorname{Top}_c\!\left(\bigcup_j
                    \operatorname{Top}_c(U_j)\right).
 \label{eq:5-topc-union}
\end{equation}
The statement includes universes smaller than $c$.
\end{proposition}
\begin{proof}
Every global winner is in $V$, and every other member of $V$ follows
the global winners. Selection from $V$ therefore recovers them.
For the independent-group special case, let $x$ belong to the global
top-$c$ and to group $U_j$. Fewer than $c$
candidates precede $x$ globally, so fewer than $c$ precede it within $U_j$.
Consequently $x\in\operatorname{Top}_c(U_j)$, and the local union contains
every global winner. Selecting its first $c$ therefore recovers precisely
the global winners in the same order. If the whole universe has fewer
than $c$ members, every member is retained.

For merging, the first $c$ of two sorted lists are their union's
top-$c$. A discarded candidate already has at least $c$ predecessors
within that union and cannot enter the top-$c$ of any larger union.
Induction over the merge tree establishes the claim.
\end{proof}
\end{resultblock}

\begin{figure}[tbp]
\centering
% Schematic only: illustrative keys, not experimental observations.
\begin{tikzpicture}[
  font=\small\linespread{1}\selectfont,
  box/.style={draw=ThesisNavy!65,rounded corners=2pt,align=center,
    inner sep=5pt,text width=2.65cm,minimum height=2.05cm},
  merged/.style={box,text width=5.55cm,minimum height=1.05cm,
    fill=ThesisNavy!5},
  arr/.style={-{Latex[length=2mm]},draw=ThesisNavy!80,thick}]
\node[box] (g0) at (0,0) {$\operatorname{Top}_3(U_0)$\\[2pt]
  $\mathbf{(1,8)}$\\$\mathbf{(2,6)}$\\$(2,9)$};
\node[box] (g1) at (3.5,0) {$\operatorname{Top}_3(U_1)$\\[2pt]
  $(4,1)$\\$(5,4)$\\\textit{length 2}};
\node[box] (g2) at (7,0) {$\operatorname{Top}_3(U_2)$\\[2pt]
  $\varnothing$\\\textit{length 0}};
\node[box] (g3) at (10.5,0) {$\operatorname{Top}_3(U_3)$\\[2pt]
  $\mathbf{(2,3)}$\\$(7,2)$\\$(8,5)$};
\node[merged] (m0) at (1.75,-2.5) {Merge and retain first 3\\[2pt]
  $\mathbf{(1,8)},\ \mathbf{(2,6)},\ (2,9)$};
\node[merged] (m1) at (8.75,-2.5) {Merge and retain first 3\\[2pt]
  $\mathbf{(2,3)},\ (7,2),\ (8,5)$};
\node[merged,text width=6.8cm,fill=ThesisNavy!12] (out) at (5.25,-4.3)
  {Global primary top-3\\[2pt]
   $\mathbf{(1,8)},\ \mathbf{(2,3)},\ \mathbf{(2,6)}$};
\draw[arr] (g0.south) -- (m0.north west);
\draw[arr] (g1.south) -- (m0.north east);
\draw[arr] (g2.south) -- (m1.north west);
\draw[arr] (g3.south) -- (m1.north east);
\draw[arr] (m0.south) -- (out.north west);
\draw[arr] (m1.south) -- (out.north east);
\end{tikzpicture}

\caption[Local coverage and truncated pairwise merging]{Local coverage and truncated pairwise merging, illustrated with
$c=3$ and $\gscan=4$. Entries are illustrative $(\primarydist,p)$ keys; bold ones are global winners. Short/empty groups contribute their valid prefixes. Merges keep at most $c$ in the same order.}
\label{fig:5-global-selection}
\end{figure}

The CUDA implementation maintains a shared-memory buffer of $C_{\rm buf}$ key
pairs, where $C_{\rm buf}$ is a power of two and $C_{\rm buf}\geq c+T_{\rm blk}$ for a scan block of
$T_{\rm blk}$ threads. Bitonic compaction keeps up to $c$ valid keys and sets the threshold to the last key once $c$ exist. Before that point the
threshold is a sentinel above every valid key. Between compactions, admission requires $K\prec\lambda_j$. The threshold has $c$ witnesses, so remains an upper bound on the prefix $c$th key as candidates arrive.

Compaction leaves at most $c$ entries; a chunk adds at most $T_{\rm blk}$:
\[
 |B|+|\Delta V|\leq c+T_{\rm blk}\leq C_{\rm buf}.
\]
This supplies \cref{lem:5-buffer}'s capacity premise without thread quotas. With \cref{prop:5-shared-bound,prop:5-local-coverage}, it preserves global winners through filtering and merging.

Materialising scores costs $O(n_c)$ writes/storage. Fused scoring compacts local buffers and emits $O(\gscan c)$ pairs. Register selection is another option \cite{johnson_gpu}.

\subsection{Merge implementation, storage and boundary cases}

At $\gscan>1$, a 256-thread block merges each query over $\log_2\gscan$ rounds. Each pair emits $\min(c,a+b)$ entries using Merge Path slices \cite{gpu_merge_path}.
All comparisons use \cref{eq:5-primary-key}; block barriers separate
rounds. \Cref{fig:5-global-selection} illustrates the truncation invariant.
The tree emits at most $c(\gscan-1)$ key pairs across its internal nodes,
in addition to the partition searches and synchronisation work.

The two global score/position workspaces have capacities
$[b_{\rm cap},\gscan,c]$ and $[b_{\rm cap},\gscan/2,c]$ and alternate between
rounds. With four-byte scores and positions, their combined payload is
\begin{equation}
 S_{\rm pairs}=8b_{\rm cap}c\left(\gscan+\frac{\gscan}{2}\right)
 \quad\text{bytes},\qquad \gscan>1.
 \label{eq:5-group-workspace}
\end{equation}
Counts add $4b_{\rm cap}\gscan$ bytes and provisioned bounds add $8b_{\rm cap}$, beyond survivor/other workspace. Allocation grows before capture and stays reserved for reuse. Merging stages lengths, not all keys, in shared memory.

Valid counts exclude padding:
\[
 K_{\rm sort\ pad}=(+\infty,\mathrm{INT\_MAX}),\qquad
 K_{\rm output\ pad}=(+\infty,-1).
\]
Valid positions are below $\mathrm{INT\_MAX}$; empty groups have count zero. Refinement skips negative positions, and merges exclude them. At $\gscan=1$, survivors go directly to output.

$\gscan=0$ uses one block per query with distance-exact selection and free cutoff-tie IDs. Lossy \texttt{K\_LOCAL=4} requires \texttt{JHQ\_EXACT\_TOPCK=0}; paper-faithful builds disable it.

For nonempty scanned set $V$, let $h=\min(c,|V|)$ and $d_h$ be its $h$th score. Distance-only compaction yields
\[
 \{x:d_x<d_h\}\subseteq B\subseteq\{x:d_x\leq d_h\},\qquad |B|=h.
\]
A rejection has $c$ witnesses with no larger score; only tied identities remain free. T4 checks this contract and grouped full-key order by CPU sorting GPU-produced scores (\cref{sec:7-strict-selection}).